\section*{参考文献}
\sloppy
\begin{thebibliography}{99}\small
\bibitem{cumcm}全国大学生数学建模竞赛官网，\url{https://www.mcm.edu.cn/}。
\bibitem{welzl}E. Welzl, Smallest enclosing disks (balls and ellipsoids), LNCS 555, 1991, 359--370.
\bibitem{toussaint}G. T. Toussaint, Solving geometric problems with the rotating calipers, IEEE MELECON, 1983.
\bibitem{tokekar2012}M. Tokekar, N. Karnofsky, and V. Isler, Sensor placement for a bearing-only target localization problem, ICRA, 2012, DOI: \url{https://doi.org/10.1109/ICRA.2012.6225244}.
\bibitem{tokekar2013}M. Tokekar and V. Isler, Robust 2-D target localization with bounded bearing errors, ICRA, 2013, DOI: \url{https://doi.org/10.1109/ICRA.2013.6630920}.
\bibitem{dogancay}U. Doğançay, On the optimality of target localization from bearing-only measurements, IEEE TAES, 2012, DOI: \url{https://doi.org/10.1109/TAES.2012.6178054}.
\bibitem{song}J. Song et al., Finding transient radio sources with mobile robots, IEEE TRO, 2012, DOI: \url{https://doi.org/10.1109/TRO.2012.2183069}.
\bibitem{fields2cover}Fields2Cover: an open-source coverage path planning library, \url{https://github.com/Fields2Cover/Fields2Cover}。
\bibitem{sisl}SISL InformativePathPlanning, \url{https://github.com/sisl/InformativePathPlanning}。
\bibitem{coverage}GradualScholar, Coverage Mission Planning, \url{https://github.com/GradualScholar/Coverage_Mission_Plannin.git}。
\bibitem{cmu}Carnegie Mellon University, Active range and bearing based radiation source localization, \url{https://publications.ri.cmu.edu/active-range-and-bearing-based-radiation-source-localization}。
\bibitem{award2020}同济大学团队，2020 CUMCM B《沙漠游戏》（公开仓库自述为全国一等奖），\url{https://github.com/seanys/CUMCM2020-Desert-Game/blob/main/沙漠游戏.pdf}。
\bibitem{award2022}浙江大学团队，2022 CUMCM A《波浪能最大输出功率优化设计》（公开仓库自述为 A 题国一），\url{https://github.com/zjuyfy/CUMCM/blob/master/AProblem.pdf}。
\bibitem{award2023}2023 CUMCM A《定日镜场优化设计模型》（公开仓库自述为国家级一等奖），\url{https://github.com/linggm3/2023_CUMCM_National-First-Prize/blob/main/定日镜场优化设计模型.pdf}。
\bibitem{award2024}2024 CUMCM A《板凳龙运动状态及路线研究》（公开仓库自述为全国一等奖），\url{https://github.com/cny123222/CUMCM-2024A-Bench-Dragon/blob/main/OurPaper.pdf}。
\bibitem{award2025}2025 CUMCM B《红外干涉法测量半导体外延层厚度》（公开仓库自述为 National First Prize），\url{https://github.com/CUMCM-2025B-Team/CUMCM-2025-Problem-B/blob/main/25国赛.pdf}。
\end{thebibliography}

\paragraph{国奖论文写作样本。}本轮另行整理了至少五篇可核验的国赛获奖论文，比较其“摘要—问题分析—模型—算法—检验—评价”的共同结构、短句化结果表达和表格化论证方式，详见仓库文档 \texttt{docs/award\_paper\_writing\_review.md}。正式提交前仅保留与本文方法直接相关、且能核验来源的引用。
